\documentclass{article}
\usepackage[utf8]{inputenc}
\usepackage[T1]{fontenc}
\usepackage{geometry}
\geometry{
    a4paper,
    left=2.5cm,
    right=2.5cm,
    top=2.5cm,
    bottom=2.5cm
}
\usepackage{graphicx}
\usepackage{float}
\usepackage{hyperref}
\usepackage[english]{babel}
\usepackage{amsmath}
\usepackage{amssymb}
\usepackage{array}
\usepackage{booktabs}
\usepackage{caption}
\usepackage{subcaption}
\usepackage{enumitem}
\usepackage{fancyhdr}
\usepackage{titlesec}
\usepackage{url}
\usepackage{textcomp}
\usepackage{xcolor}
\usepackage{listings}
\usepackage{multirow}
\usepackage{longtable}
\usepackage{tabularx}
\usepackage{times}

% Set up page style
\pagestyle{fancy}
\fancyhf{}
\rfoot{\thepage}
\renewcommand{\headrulewidth}{0pt}

% Title formatting
\titleformat{\chapter}[display]{\normalfont\bfseries}{\chaptertitle\thechapter}{0pt}{\Huge}
\titlespacing*{\chapter}{0pt}{50pt}{40pt}

% Document info
\title{Template Document}
\author{Academic Support Agent}
\date{\today}

\begin{document}
\maketitle

\thispagestyle{title}



% Problem 1
\problem Show that there exists no nontrivial unramified extensions of $\Q$. \pspace

\solution  If $K/\Q$ is a nontrivial number field, then $|\disc K|>1$. But then $\disc K$ has a prime factor so that some prime ramifies in $K$. \qed \pspace



% Problem 2
\problem Complete the following:
        \begin{enumerate}[(a)]
	\item How does one prove a cotheorem?
	\item Compute $\ds \int \cos x \;dx$.
	\item How does one square $\twomatrix{a}{b}{c}{d}$?       
        \end{enumerate}

\solution
\begin{enumerate}[(a)]
\item Use rollaries.

\item We have
	\begin{equation} \label{eq:integral}
	\int \cos x \;dx= \sin x + C	
	\end{equation}
We can check \eqref{eq:integral}:
	\[
	\dfrac{d}{dx} \left( \sin x + C \right)= \cos x
	\]

\item This is routine. 
\end{enumerate} \qed \pspace



% Problem 3
\problem Prove that $\sqrt{2}$ is irrational. \pspace

\pf Assume that $\sqrt{2}= \dfrac{a}{b}$, where $a,b \in \Z$. Without loss of generality, we may assume $\gcd(a,b)= 1$. Then we have
	\begin{align}
	\sqrt{2}&= \dfrac{a}{b} \nonumber \\
	\sqrt{2}^2&= \left( \dfrac{a}{b} \right)^2 \label{eq:implication1} \\
	2&= \dfrac{a^2}{b^2} \nonumber \\
	a^2&= 2b^2 \label{eq:implication2}
	\end{align}
But then from \eqref{eq:implication2}, we know that $a^2$ is even so that $a$ is even. But then we must have
	\[
	2a^2= b^2
	\]
so that $b^2$ is even, implying $b$ is even. But then $\gcd(a,b) \geq 2$, a contradiction. \qed \\
\end{document}